% !TeX program = xelatex
\documentclass[aspectratio=169]{nextbeamer_zh}

\UseTemplateSet{
  layout = beamer,
  globalfonts = {cmu,shanggu},
  features = {tables,image}
}

\title{中文 Beamer 模板示例}
\author{作者姓名}
\date{\today}

\begin{document}

\begin{frame}
  \titlepage
\end{frame}

\begin{frame}{目錄}
  \tableofcontents
\end{frame}

\section{研究背景}

\begin{frame}{研究背景}
  \begin{itemize}
    \item 在這一頁簡述題目背景與研究動機。
    \item 可以直接替換成課堂報告、論文發表或讀書報告內容。
    \item 中英文混排與專有名詞可直接保留在正文中。
  \end{itemize}
\end{frame}

\section{研究問題}

\begin{frame}{研究問題}
  \begin{enumerate}
    \item 研究對象是什麼？
    \item 現有研究如何處理？
    \item 你的觀察或方法在哪裡？
  \end{enumerate}
\end{frame}

\section{材料示例}

\begin{frame}{表格示例}
  \begin{tabular}{ll}
    項目 & 說明 \\
    \hline
    材料 & 文本樣本 \\
    方法 & 分析與比較 \\
    結果 & 初步觀察 \\
  \end{tabular}
\end{frame}

\begin{frame}{圖片佔位}
  \centering
  \rule{0.75\linewidth}{0.35\linewidth}
\end{frame}

\section{結語}

\begin{frame}{結語}
  \begin{itemize}
    \item 這份模板不是最小示例，而是可以直接開始製作投影片的骨架。
    \item 後續可在此加入更多 section、圖片、表格與引文頁面。
  \end{itemize}
\end{frame}

\end{document}
